% Interfaz computable CG-008; formalización nueva de un lector residual.
% Productor: los caracteres e intervalos regionales construidos previamente.
% No identifica la componente numérica con el estado enriquecido completo.

\section{Initialisation résiduelle calculable et compatibilité des émissions régionales}
\label{sec:residuo-regional-computable}

La réalisation numérique d'une section régionale admet une description
exécutable par intervalles rationnels. Cette description permet d'initialiser
le résidu, de le mettre à jour et de récupérer ses blocs sans introduire une chaîne
décimale préalablement fixée. Son origine est le producteur régional du
\cref{x_reciprocidad:sec:generacion} : le caractère,
son orientation et sa normalisation sont d'abord déterminés ; ses approximations
rationnelles sont ensuite construites. Le présent développement formalise la composante résiduelle de cette
lecture. La trajectoire, le feuillet, les incidences et la mémoire de l'état
enrichi conservent simultanément leurs opérations propres.

L'ordre des opérations est donc
\[
 \begin{gathered}
 \text{état régional généré}
 \longrightarrow \text{intervalles rationnels du caractère}
 \\
 \longrightarrow \text{nom rationnel du résidu}
 \longrightarrow \text{émissions positionnelles compatibles}.
 \end{gathered}
\]
La comparaison avec la réalisation archimédienne apparaît dans la démonstration
de correction. L'algorithme reçoit des intervalles ; aucune de ses instructions
ne reçoit un nombre réel conventionnel ni une liste de chiffres cibles.

\subsection{Nom rationnel produit par un lecteur régional}
\label{subsec:residual-nombre-racional}

On fixe une région et l'on note \(j\mapsto[\ell_j,u_j]\) son producteur
rationnel. Les conditions de l'interface sont
\begin{equation}
 \ell_j\leq v\leq u_j,\qquad
 \lim_{j\to\infty}(u_j-\ell_j)=0,
 \qquad \ell_j,u_j\in\mathbb Q.
 \label{eq:residual-contrato-productor}
\end{equation}
Ici, \(v\) désigne la réalisation exacte du caractère pour formuler la
correction ; la procédure qui calcule \(\ell_j,u_j\) provient du lecteur
régional. Les intervalles de clôture, de propagation et d'autoéchelle déjà
construits satisfont ces conditions. Dans l'implémentation, son interface
est une fonction qui, étant donné \(j\), renvoie deux fractions exactes.

L'emboîtement est obtenu, si nécessaire, au moyen de l'opération finie
\begin{equation}
 L_j=\max_{0\leq i\leq j}\ell_i,
 \qquad U_j=\min_{0\leq i\leq j}u_i,
 \qquad I_j=[L_j,U_j].
 \label{eq:residual-anidacion}
\end{equation}

\begin{proposition}[Emboîtement et conservation de la réalisation]
\label{prop:residual-anidacion}
Les intervalles \(I_j\) sont rationnels, non vides et emboîtés. Ils contiennent
\(v\), leur largeur tend vers zéro et \(\bigcap_j I_j=\{v\}\).
\end{proposition}
\begin{proof}
Chaque borne inférieure est inférieure ou égale à \(v\), et chaque borne supérieure est
supérieure ou égale à \(v\) ; la même propriété vaut pour le maximum et le minimum
finis. Lorsque \(j\) augmente, le maximum croît et le minimum décroît.
De plus, \(0\leq U_j-L_j\leq u_j-\ell_j\). Si deux points appartenaient à
l'intersection, leur distance serait bornée par chaque largeur et serait nulle.
L'appartenance de \(v\) atteste la non-vacuité.
\end{proof}

\subsection{Préfixe initial et représentation du résidu}
\label{subsec:residual-inicializacion}

Soient \(b\geq2\) une base entière et \(n_0\geq0\) une profondeur.
On suppose que la réalisation évite les frontières correspondantes :
\begin{equation}
 b^n v\notin\mathbb Z\qquad(n\geq0).
 \label{eq:residual-no-frontera}
\end{equation}
L'irrationalité des trois valeurs régionales implique cette propriété ;
sa démonstration ou son emploi comme hypothèse relève du lecteur fixé.
Pour une profondeur concrète, l'évitement aux échelles utilisées
suffit. Les réalisations à développement terminal exigent, en outre,
un décideur exact de frontière et une convention de représentation.

Nous définissons le prédicat rationnel
\[
 \mathsf{Sep}(j,n):\quad
 \lfloor b^n L_j\rfloor=\lceil b^n U_j\rceil-1.
\]
Il est décidable au moyen d'opérations entières sur les numérateurs et les
dénominateurs. Soient
\begin{equation}
 j_0=\min\{j:\mathsf{Sep}(j,n_0)\},
 \qquad N_0=\lfloor b^{n_0}L_{j_0}\rfloor.
 \label{eq:residual-prefijo-inicial}
\end{equation}
L'entier \(N_0\) est calculé à partir du producteur. Son emploi comme antécédent
de la représentation résiduelle n'exige pas de l'introduire comme donnée littérale.

\begin{theorem}[Initialisation calculable du résidu régional]
\label{thm:residual-inicializacion}
L'indice \(j_0\) existe. Pour \(j\geq j_0\), les intervalles
\begin{equation}
 E_{0,j}=
 \bigl[b^{n_0}L_j-N_0,\ b^{n_0}U_j-N_0\bigr]
 \label{eq:residual-nombre-inicial}
\end{equation}
sont non vides, emboîtés et contenus dans \([0,1]\). Leur largeur tend vers
zéro et leur intersection consiste en un unique élément
\(\varepsilon_0\in(0,1)\). Dans la réalisation archimédienne, on a
\begin{equation}
 N_0=\lfloor b^{n_0}v\rfloor,
 \qquad b^{n_0}v=N_0+\varepsilon_0.
 \label{eq:residual-identificacion-inicial}
\end{equation}
La suite rationnelle \((E_{0,j})_{j\geq j_0}\) constitue une
représentation calculable du résidu.
\end{theorem}
\begin{proof}
Le point \(b^{n_0}v\) est à l'intérieur d'une cellule d'extrémités
entières. Sa distance à ces extrémités est positive. Lorsque la largeur de
\(b^{n_0}I_j\) est inférieure à cette distance, les deux extrémités appartiennent à la
même cellule et \(\mathsf{Sep}(j,n_0)\) est vérifié. Il existe donc un
premier indice de séparation. À tout indice séparé, l'entier
\(\lfloor b^{n_0}v\rfloor\) se trouve entre les deux entiers égaux
du prédicat et coïncide avec eux. L'égalité initiale situe les extrémités
de \(b^{n_0}I_{j_0}\) entre \(N_0\) et \(N_0+1\). L'emboîtement
préserve cette inclusion. La transformation affine de
\eqref{eq:residual-nombre-inicial} conserve l'emboîtement et la non-vacuité et
multiplie les largeurs par \(b^{n_0}\). Leur intersection est le point de
\eqref{eq:residual-identificacion-inicial} ; l'évitement de frontière le
situe dans \((0,1)\). Chaque extrémité est calculée au moyen de fractions exactes,
et la recherche se termine par l'argument précédent.
\end{proof}

L'identité \eqref{eq:residual-identificacion-inicial} est une conclusion
de correction. La représentation exécutable est la suite
\eqref{eq:residual-nombre-inicial}, construite à partir d'intervalles
régionaux. Cette distinction sépare la production d'un nom rationnel de la
reconnaissance de sa valeur.

\subsection{Mise à jour résiduelle et émission par arithmétique exacte}
\label{subsec:residual-actualizacion}

Supposons construit l'état de lecture \((n_0+k,N_k,j_k)\), ainsi
que le nom
\begin{equation}
 E_{k,j}=
 \bigl[b^{n_0+k}L_j-N_k,\ b^{n_0+k}U_j-N_k\bigr],
 \qquad j\geq j_k.
 \label{eq:residual-nombre-k}
\end{equation}
On augmente \(j\) à partir de \(j_k\) jusqu'à ce que
\[
 \left\lfloor b\inf E_{k,j}\right\rfloor
 =\left\lceil b\sup E_{k,j}\right\rceil-1.
\]
Le premier indice qui satisfait cette égalité est noté \(j_{k+1}\), et
l'entier commun est noté \(d_k\). La mise à jour est
\begin{equation}
 N_{k+1}=bN_k+d_k,
 \qquad E_{k+1,j}=bE_{k,j}-d_k
 \quad(j\geq j_{k+1}).
 \label{eq:residual-actualizacion-exacta}
\end{equation}
La multiplication d'un intervalle par \(b>0\) et la soustraction de \(d_k\)
agissent sur ses deux extrémités. On conserve l'indice de lecture atteint,
la profondeur et le bloc émis.

\Needspace{20\baselineskip}
\begin{theorem}[Terminaison à profondeur arbitraire et commutation positionnelle]
\label{thm:residual-compatibilidad-universal}
Sous \eqref{eq:residual-contrato-productor} et
\eqref{eq:residual-no-frontera}, chaque mise à jour se termine et, pour tout
\(k\geq0\), on a
\begin{align}
 N_k&=\lfloor b^{n_0+k}v\rfloor,
 &0&\leq d_k<b,\label{eq:residual-prefijos-k}\\
 \varepsilon_{k+1}&=b\varepsilon_k-d_k,
 &d_k&=\lfloor b\varepsilon_k\rfloor,\label{eq:residual-recurrencia}\\
 b^{n_0+k}v&=N_k+\varepsilon_k,
 &0&<\varepsilon_k<1.\label{eq:residual-invariante}
\end{align}
Si \(P(b,n)\) est l'émission par séparation rationnelle du
\cref{thm:df-emision-correcta}, alors
\begin{equation}
 N_k=P(b,n_0+k),\qquad
 d_k=P(b,n_0+k+1)\bmod b.
 \label{eq:residual-igualdad-publicacion}
\end{equation}
Les indices d'arrêt des deux implémentations peuvent être distincts ;
leurs entiers émis coïncident exactement.
\end{theorem}
\begin{proof}
L'initialisation fournit les assertions pour \(k=0\). En les supposant
au niveau \(k\), la multiplication de l'invariant par \(b\) donne
\[
 b^{n_0+k+1}v=bN_k+b\varepsilon_k.
\]
L'évitement de frontière implique \(b\varepsilon_k\notin\mathbb Z\).
Les intervalles \(bE_{k,j}\) sont emboîtés, contiennent ce point et ont une
largeur tendant vers zéro. L'argument de séparation du théorème précédent
garantit la terminaison et fixe
\(d_k=\lfloor b\varepsilon_k\rfloor\). Comme
\(0<\varepsilon_k<1\), le bloc appartient à
\(\{0,\ldots,b-1\}\). Soustraire \(d_k\) transforme les intervalles en
un nom emboîté du résidu suivant et conserve son appartenance à
\((0,1)\). La partie entière de l'égalité précédente est
\(bN_k+d_k\), ce qui prouve le nouvel invariant. La récurrence vaut pour
chaque profondeur finie, sans imposer de profondeur maximale.
Enfin, le théorème d'émission rationnelle identifie également
\(P(b,n)\) à \(\lfloor b^n v\rfloor\). L'égalité des
préfixes et la formule du résidu modulo \(b\) démontrent
\eqref{eq:residual-igualdad-publicacion}.
\end{proof}

\begin{corollary}[Troncature, récupération et bornes résiduelles]
\label{cor:residual-recuperacion}
Pour \(k,h\geq0\),
\[
 \left\lfloor\frac{N_{k+h}}{b^h}\right\rfloor=N_k,
 \qquad N_k=\frac{N_{k+1}-d_k}{b},
 \qquad \varepsilon_k=\frac{\varepsilon_{k+1}+d_k}{b}.
\]
La largeur du nom résiduel à l'indice \(j\) est
\(b^{n_0+k}(U_j-L_j)\). Par conséquent, toute précision rationnelle
positive pour le résidu est atteinte après un nombre fini de consultations du
producteur. La profondeur et les blocs archivés déterminent la
récupération dans la composante de lecture.
\end{corollary}
\begin{proof}
La division du préfixe étendu récupère le préfixe précédent par
\eqref{eq:residual-prefijos-k}. Les deux formules de récupération sont
obtenues par résolution dans \eqref{eq:residual-actualizacion-exacta} et
\eqref{eq:residual-recurrencia}. La formule de largeur découle immédiatement de
\eqref{eq:residual-nombre-k} ; sa convergence vers zéro fournit la
terminaison de la recherche de précision.
\end{proof}

\subsection{Correspondance des profondeurs ternaires et décimales}
\label{subsec:residual-bases}

Dans la normalisation fractionnaire, on écrit \(v=m+r\), où
\(m=P(b,0)\) est obtenu du même producteur. Le préfixe fractionnaire est
\(C_n=P(b,n)-mb^n\). Par calcul entier,
\[
 b^n v-P(b,n)=b^n r-C_n.
\]
Cette égalité réconcilie le préfixe qui inclut la partie entière avec celui qui
commence à zéro, en conservant un résidu identique. La soustraction de \(m\)
est une transformation ultérieure du nom régional.

Le regroupement tritique utilise \(b=729=3^6\). Six blocs correspondent à
\(36\) trits ; la spécialisation \(n_0=6\) fournit un nom
rationnel du résidu de cette profondeur. Le regroupement décimal utilise
\(b=1000=10^3\) ; six blocs correspondent à \(18\) chiffres décimaux.
Ce sont des coordonnées de lecture différentes du même caractère. Leur compatibilité
s'exprime par les cylindres qui contiennent cette réalisation et par les
identités de changement d'échelle ; elle n'exige pas d'égaler leurs résidus numériques.

L'antécédent résiduel de
\cref{sec:c27-generacion-coinductiva-profundidad-arbitraria} distingue
l'état interne de sa publication. La présente interface apporte une
représentation rationnelle exécutable de la composante numérique, avec la même
récurrence résiduelle une fois sa lecture fixée. Pour l'identifier à une autre
réalisation résiduelle de cette section, il suffit d'établir que les deux lisent le même
caractère, à la même profondeur et avec la même normalisation : l'identité
\(b^n v=N_n+\varepsilon_n\) et l'unicité du préfixe imposent alors
l'égalité des résidus. Cette comparaison est une égalité de lecteurs.
L'évolution des positions, des feuillets, de l'incidence et de la mémoire demeure régie
par le TPK enrichi et par leurs théorèmes de transport respectifs.

\subsection{Contrat de l'implémentation exacte}
\label{subsec:residual-implementacion}

L'implémentation de cette interface emploie exclusivement des fractions
rationnelles, des puissances entières, le maximum, le minimum, la division avec plancher et la division
avec plafond. Le producteur d'intervalles est un argument explicite du
programme. Une consultation ne renvoie pas \(v\) : elle renvoie ses deux bornes engendrées.
L'état d'exécution conserve la base, la profondeur, le préfixe, l'indice atteint
et le registre ordonné des émissions. Pour chaque étape, le registre garde le
nom résiduel précédent à l'indice décisif, le bloc et le nom
transformé ; les deux récupérations du corollaire sont vérifiées au moyen
d'égalités rationnelles.

Les preuves universelles précédentes justifient la terminaison et la
compatibilité sous le contrat du producteur. Les essais finis vérifient
l'implémentation de ces opérations et de leurs identités ; la généralité
du résultat réside dans la récurrence et dans la séparation rationnelle des cellules.
Cette formalisation rassemble une interface calculable de lecture et de récupération
sans remplacer la sélection régionale ni attribuer à la composante numérique
les opérations supplémentaires de l'état enrichi.
